\subsection{Структура программы}

Было принято решение, что при разработке данного программного средства будут использоваться пять основных модулей и девять подмодулей:
\begin{itemize}
    \item Utils -- вспомогательные функции для работы с путями и форматированием данных;
    \item Settings -- управление настройками приложения, темами, шрифтами и их сохранение в реестре Windows;
    \item FileSystem -- операции с файловой системой (навигация по директориям, работа с деревом каталогов, получение информации о файлах);
    \item Network -- сканирование сетевых устройств, перечисление сетевых ресурсов, работа с сетевыми дисками и USB устройствами;
    \item UI -- графический интерфейс приложения, включающий следующие подмодули:
    \begin{itemize}
        \item List -- управление отображением файлов в различных режимах (подробно, крупные значки, мелкие значки);
        \item Layout -- размещение элементов интерфейса и адаптация к изменению размера окна;
        \item Navigation -- навигация по файловой системе, управление историей переходов, обновление состояния интерфейса;
        \item Menu -- создание и управление меню приложения, обработка горячих клавиш;
        \item Theme -- применение тем оформления (светлая, тёмная, системная) и настройка шрифтов;
        \item SubClass -- отрисовка элементов управления для поддержки тёмной темы;
        \item Window -- хранение и управление состоянием окна, инкапсуляция доступа к данным;
        \item Control -- создание и настройка всех элементов интерфейса;
        \item WindowProc -- обработка всех сообщений Windows и координация работы модулей;
    \end{itemize}
\end{itemize}

Каждый модуль будет содержать интерфейсный файл для доступа к функциональности модуля из других модулей.